\subsection{Milestone 1 -- Costruzione del dataset}

\subsubsection{Raccolta e filtraggio delle versioni}

Le versioni del progetto sono state scaricate tramite le API REST di Jira. Sono state
considerate esclusivamente le versioni di tipo \emph{release}, ovvero quelle per cui il
campo \texttt{released} risulta vero e che dispongono di una data di rilascio
(\texttt{releaseDate}). Le versioni prive di uno di questi due attributi sono state
scartate. In totale, per Apache OpenJPA, sono state identificate 42 versioni valide.

I ticket estratti dalla API REST di Jira sono stati filtrati attraverso la seguente jql:
\textit{"issuetype = Bug AND status in (Closed, Resolved) AND resolution = Fixed"}

Per la costruzione degli snapshot e il calcolo delle metriche è stato deciso di
considerare solo il \textbf{primo 33\% delle release} in ordine cronologico, come da
specifica del progetto. Questo equivale a 14 release su 42. Questa scelta è stata fatta in vista di ridurre al minimo il numero di bug dormienti e/o "snoring" del progetto analizzato, i quali compaiono statisticamente sempre piu spesso avanzando verso le versioni recenti.

\subsubsection{Raccolta dei ticket Jira}

Ogni ticket è modellato con i seguenti campi rilevanti:
\begin{itemize}
    \item \textbf{created}: data di creazione del ticket, in formato \texttt{yyyy-MM-dd}.
    \item \textbf{affectedVersions (AV)}: versioni nelle quali il sistema è affetto dal bug.
    \item \textbf{fixVersions (FV)}: versioni nelle quali il bug è stato corretto.
    \item \textbf{openingVersion (OV)}: prima versione del progetto cronologicamente
    precedente alla data di creazione del ticket.
    \item \textbf{injectedVersion (IV)}: la AV meno recente del ticket, calcolata
    dinamicamente in quanto fortemente accoppiata alla lista di AV, modificata più volte
    durante i check di consistenza.
\end{itemize}

\subsubsection{Arricchimento dei ticket}

Prima dei check di consistenza, i ticket vengono arricchiti tramite le seguenti operazioni:

\begin{enumerate}
    \item \textbf{Calcolo della Opening Version}: la OV è la prima versione del progetto
    precedente cronologicamente alla data di creazione del ticket. I ticket privi di data
    di creazione o con data antecedente a qualsiasi versione del progetto vengono scartati. Questo poichè non è ragionevole lavorare con ticket privi di informazioni basilari o che sono stati aperti prima ancora di aver avuto la prima release di un progetto, i quali avrebbero bisogno di un approccio dedicato, senza il quale rischieremmo di portare rumore nei dati. Questo approccio sceglie la qualità dei dati rispetto alla quantità.

    \item \textbf{Riduzione a una singola Fix Version}: per i ticket con più FV viene
    mantenuta solo la più recente, ritenuta la più affidabile. Anche in questo caso è stato preferito lavorare con informazioni piu precise al costo di perderne una parte minima.

    \item \textbf{Rimozione delle versioni senza data}: le versioni referenziate nei campi
    dei ticket prive di data di rilascio vengono rimosse direttamente dai campi del ticket. Questo sulla base della scelta di non usare versioni senza release date.

    \item \textbf{Impostazione delle AV a partire dalle FV}: se un ticket non possiede AV
    ma ha FV, le FV tranne l'ultima vengono interpretate come AV. Questa scelta potrebbe essere
    preferibile all'utilizzo di Proportion, che per definizione costituisce una stima. L'interpretazione fatta mira a raccogliere informazioni dai ticket sulla base di un'ipotesi, che potrebbe essere corretta o meno.

    \item \textbf{Rimozione delle versioni non più valide}: dopo i check di consistenza,
    le versioni rimosse vengono eliminate anche dai riferimenti presenti nei campi dei
    ticket.
\end{enumerate}

\subsubsection{Check di consistenza}

I check di consistenza sono componenti modulari e indipendenti, ciascuno con una singola
responsabilità. Alcuni check sui ticket devono essere eseguiti solo dopo determinate
operazioni di enriching, motivo per cui vengono invocati in uno specifico ordine.

\paragraph{Check sui ticket}
\begin{itemize}
    \item \textbf{IssueHasKeyCheck}: il ticket deve avere una chiave identificativa.
    \item \textbf{IssueHasCreatedDateCheck}: il ticket deve avere una data di creazione.
    \item \textbf{IssueHasFixVersionCheck}: il ticket deve avere almeno una FV.
    \item \textbf{IssueHasAffectedVersionCheck}: il ticket deve avere almeno una AV.
    \item \textbf{IssueFixVersionHasReleaseDateCheck}: la FV deve avere una data di rilascio.
    \item \textbf{IssueAllVersionsHaveReleaseDateCheck}: tutte le versioni referenziate
    devono avere una data di rilascio.
    \item \textbf{IssueCreatedAfterOldestVersionCheck}: la data di creazione deve essere
    successiva alla data della versione più vecchia del progetto.
    \item \textbf{IssueFixVersionAfterAffectedVersionCheck}: la FV deve essere
    necessariamente successiva a ogni AV. I ticket non conformi vengono scartati. Si è
    preferito scartare il ticket piuttosto che rimuovere le AV inconsistenti, per evitare
    di modificare soggettivamente il significato del ticket.
    \item \textbf{IssueFixVersionAfterEqualOpeningVersionCheck}: deve esistere almeno una
    FV con $\text{FV} \geq \text{OV}$. Eseguito solo su ticket con singola FV.
\end{itemize}

\paragraph{Check sulle versioni}
\begin{itemize}
    \item \textbf{VersionHasNameCheck}: la versione deve avere un nome.
    \item \textbf{VersionHasReleaseDateCheck}: la versione deve avere una data di rilascio.
    \item \textbf{VersionIsReleasedCheck}: la versione deve essere di tipo release.
\end{itemize}

\subsubsection{Stima delle Affected Versions con Proportion}

Per i ticket privi di AV dopo le operazioni di enriching, le AV vengono stimate tramite la
tecnica \textbf{Proportion} con approccio \emph{cold start}. Il valore $P$ è calcolato
come la media di $(FV - IV) / (FV - OV)$ su tutti i ticket con FV, IV e OV disponibili.
I ticket con $FV = OV$ vengono esclusi dal calcolo per evitare la divisione per zero.
La IV stimata viene calcolata come:

\[
IV = FV - P \cdot (FV - OV)
\]

Se la IV predetta coincide con la FV il ticket viene scartato. Per Apache OpenJPA, un
ticket è stato rimosso per questa ragione.

\subsubsection{Estrazione degli snapshot di classi Java}

Gli snapshot (classe al tempo dell'ultimo commit di una release) sono stati estratti dal repository Git locale tramite comandi Git invocati
programmaticamente. Per ogni release viene identificato il tag Git corrispondente;
i file vengono elencati con \texttt{git ls-tree -r --name-only <tag>} e il codice sorgente
estratto con \texttt{git show <tag>:<path>}, salvato su disco per release.

Sono considerate esclusivamente le classi che soddisfano i seguenti criteri:
\begin{itemize}
    \item il percorso termina con \texttt{.java};
    \item il percorso non contiene directory di test (\texttt{/test/} o \texttt{/tests/});
    \item il nome del file non termina con \texttt{Test.java} o \texttt{Tests.java}.
\end{itemize}

\paragraph{Corrispondenza versioni Jira -- tag Git}

La risoluzione del tag Git per ogni versione Jira segue questa strategia:
\begin{enumerate}
    \item \textbf{Match esatto}: nome versione Jira $=$ nome tag (es.\ \texttt{1.0.0}).
    \item \textbf{Match con suffisso}: il tag inizia con il nome versione seguito da
    trattino, escludendo i tag \texttt{SNAPSHOT} (es.\ \texttt{0.9.6-incubating}).
    \item \textbf{Nessun match}: la versione viene saltata.
\end{enumerate}

Tre versioni Jira (\texttt{0.9.0}, \texttt{2.0.0-M1}, \texttt{2.0.0-M2}) non dispongono
di un tag Git e sono state escluse dagli snapshot, mantenendo però i ticket che le
referenziano. Si è scelto di non compensare le release saltate nel calcolo del 33\%, sempre per non avvicinarci troppo ai bug dormienti.

\subsubsection{Etichettatura degli snapshot (Bugginess)}

Uno snapshot (\emph{classPath}, \emph{release}) viene etichettato come \textbf{buggy} se
esiste almeno un ticket tale che:
\begin{enumerate}
    \item la release dello snapshot è inclusa tra le AV del ticket;
    \item almeno un commit il cui messaggio contiene la chiave del ticket coinvolge il
    percorso della classe.
\end{enumerate}

Il processo è basato sul seguente ragionamento: se una classe è toccata da un commit che referenzia un ticket issue, allora quella classe sarà "Buggy" in ogni AV del ticket.

Il linkage tra ticket e commit avviene tramite un indice globale costruito con un singolo
\texttt{git log --all --name-only}, che associa a ogni chiave di ticket l'insieme dei
file toccati dai commit che la citano.

\subsubsection{Calcolo delle metriche}

Per ogni snapshot vengono calcolate le seguenti metriche.

\paragraph{Metriche puntuali}
Calcolate sullo stato del file all'ultimo commit della release:
\begin{itemize}
    \item \textbf{LOC}: righe di codice non vuote.
    \item \textbf{Previous\_Release\_Code\_Smells}: violazioni PMD rilevate nella release
    immediatamente precedente.
\end{itemize}

\paragraph{Metriche cumulative}
Calcolate su tutti i commit del file dall'inizio della storia fino alla release corrente:

\begin{itemize}
    \item \textbf{LOC\_Touched}: somma di righe aggiunte e cancellate in ogni commit.
    $\text{LOC\_Touched} = \sum (\text{linesAdded} + \text{linesDeleted})$

    \item \textbf{NR}: numero totale di commit che hanno modificato il file.
    $\text{NR} = \text{count}(\text{commit sul file})$

    \item \textbf{Nfix}: numero di ticket la cui \texttt{fixVersion} corrisponde alla
    release corrente e i cui commit citano la chiave del ticket e toccano il file.\\
    $\text{Nfix} = \text{count}(\text{ticket con FV = release AND classPath} \in \text{file toccati})$

    \item \textbf{Nauth}: numero di autori distinti (per email) che hanno effettuato
    almeno un commit sul file.
    $\text{Nauth} = \text{count}(\text{authorEmail distinti})$

    \item \textbf{LOC\_Added}: somma delle sole righe aggiunte in ogni commit.
    $\text{LOC\_Added} = \sum \text{linesAdded}$

    \item \textbf{Max\_LOC\_Added}: massimo numero di righe aggiunte in un singolo commit.
    $\text{Max\_LOC\_Added} = \max(\text{linesAdded})$

    \item \textbf{Avg\_LOC\_Added}: media delle righe aggiunte per commit.
    $\text{Avg\_LOC\_Added} = \text{LOC\_Added} / \text{NR}$

    \item \textbf{Churn}: somma di righe aggiunte e cancellate per ogni commit. Misura
    il volume di riscrittura del codice nel tempo, concettualmente distinto da LOC\_Touched
    anche se numericamente equivalente.
    $\text{Churn} = \sum (\text{linesAdded} + \text{linesDeleted})$

    \item \textbf{Max\_Churn}: massimo churn registrato in un singolo commit.
    $\text{Max\_Churn} = \max(\text{linesAdded} + \text{linesDeleted})$

    \item \textbf{Avg\_Churn}: churn medio per commit.
    $\text{Avg\_Churn} = \text{Churn} / \text{NR}$

    \item \textbf{Change\_Set\_Size}: somma del numero di file committati insieme a questo
    file in ogni commit. Un valore alto indica forte accoppiamento con altre classi.
    $\text{Change\_Set\_Size} = \sum \text{changeSetSize}$

    \item \textbf{Max\_Change\_Set}: massimo numero di file committati insieme in un
    singolo commit.
    $\text{Max\_Change\_Set} = \max(\text{changeSetSize})$

    \item \textbf{Avg\_Change\_Set}: media del change set per commit.
    $\text{Avg\_Change\_Set} = \text{Change\_Set\_Size} / \text{NR}$

    \item \textbf{Age}: giorni tra il primo commit che ha introdotto il file e la data
    di rilascio della release.
    $\text{Age} = \text{releaseDate} - \text{data primo commit}$

    \item \textbf{Weighted\_Age}: Age moltiplicato per LOC\_Touched. Combina anzianità
    e volume di attività: un file vecchio ma mai modificato ha valore basso, un file
    vecchio e molto modificato ha valore alto.
    $\text{Weighted\_Age} = \text{Age} \times \text{LOC\_Touched}$
\end{itemize}

\paragraph{Metriche di commit}
Calcolate sui commit, e non sulle classi a differenza delle precedenti:
\begin{itemize}
    \item \textbf{EXP} (\textit{Developer Experience}): misura, per ciascun commit che ha interessato il file fino alla release considerata, il numero di \textit{commit-day} distinti accumulati globalmente dall'autore nell'intera repository; il valore finale è la media di tale esperienza su tutti i commit del file.
    \item \textbf{SEXP} (\textit{Subsystem Developer Experience}): applica la medesima logica restringendo il conteggio al solo sotto-sistema di appartenenza del file (identificato dal primo segmento del \textit{path}, tipicamente corrispondente al modulo Maven), catturando così la familiarità dell'autore con l'area specifica del codice.
\end{itemize}
Tali metriche sono state scelte in quanto non ridondanti rispetto a quelle già presenti (ad esempio, LA è coperta da \texttt{LOC\_Added}, NDEV da \texttt{Nauth} e NUC da \texttt{NR}) e implementabili senza modifiche architetturali rilevanti, sfruttando l'indice Git già costruito in memoria. Entrambe le metriche sono calcolate tramite ricerca binaria su liste preordinate costruite \textit{una tantum} al lancio, evitando ulteriori chiamate al processo \texttt{git}. Altre metriche quali \textit{Entropy}, NS e ND, pur di interesse, sono state escluse poiché avrebbero richiesto la definizione di un secondo indice (\textit{commit} $\to$ \textit{file}), comportando una modifica più invasiva all'architettura esistente.

\paragraph{Code smell e shift temporale}

I code smell sono rilevati tramite \textbf{PMD} con i ruleset \emph{bestpractices},
\emph{design}, \emph{errorprone} e \emph{codestyle}. La metrica è denominata \newline \textbf{Previous\_Release\_Code\_Smells} in quanto contiene il conteggio rilevato nella
release cronologicamente precedente, non in quella corrente. Questa scelta riflette
il fatto che i code smell introdotti nella release $N$ non compromettono quella release,
ma quelle successive:

\[
\text{Previous\_Release\_Code\_Smells}(\text{classe},\, N) = \text{Code\_Smells}(\text{classe},\, N-1)
\]

Gli snapshot della prima release e le classi che appaiono per la prima volta ricevono
valore 0.

\paragraph{Nota sui branch paralleli}

OpenJPA sviluppava in parallelo un branch di manutenzione \texttt{1.x} e un branch di
sviluppo \texttt{2.x}. Alcune classi introdotte nel branch \texttt{2.x} possono risultare
assenti nella release \texttt{1.x} che si interpone cronologicamente. In tali casi il
valore \texttt{Previous\_Release\_Code\_Smells} è 0, poiché la release precedente non
contiene la classe. Questo comportamento è corretto rispetto alla logica implementata;
l'alternativa di cercare l'ultima release in cui la classe era presente avrebbe introdotto
una complessità maggiore e un salto temporale implicito nella definizione di
``release precedente''.

\subsubsection{Aggiornamenti alla pipeline durante la Milestone 2}

Durante lo sviluppo della Milestone 2 sono state apportate tre modifiche ortogonali alla
pipeline della Milestone 1. Nessuna altera la logica di calcolo delle metriche o della
bugginess labeling.

\paragraph{Clone automatico del repository (package \texttt{repocloner})}

In precedenza il repository Git di OpenJPA doveva essere clonato manualmente prima
dell'esecuzione. La cartella \texttt{gitclones/openjpa} non è versionata nel repository
del progetto, quindi chi scaricava il codice doveva eseguire il clone a mano.

È stato aggiunto il package \texttt{com.isw2project.repocloner}, con due classi che
seguono il pattern service/orchestrator già usato nel resto del progetto:

\begin{itemize}
\item \texttt{RepoCloneService}: esegue il clone tramite \textbf{JGit} (libreria Java
pura, senza dipendenze dal client \texttt{git} installato sul sistema) con
\texttt{depth=1} (shallow clone) per limitare la dimensione del download.
\item \texttt{RepoCloneOrchestrator}: controlla se \texttt{<repoDir>/.git} esiste già;
se sì, salta il clone; altrimenti delega a \texttt{RepoCloneService}.
\end{itemize}

\texttt{Milestone1Main} chiama \texttt{RepoCloneOrchestrator.ensureCloned(config.getGit())}
come primo passo, prima di qualsiasi operazione Git: avviare la Milestone 1 è sufficiente
per rendere disponibile il repository localmente, senza eseguire manualmente un clone da
riga di comando. Il client \texttt{git} di sistema rimane necessario per la fase di
estrazione degli snapshot, che usa \texttt{ProcessBuilder} con \texttt{git ls-tree} e
\texttt{git show}.

Dipendenza aggiunta in \texttt{pom.xml}: \texttt{org.eclipse.jgit:org.eclipse.jgit:7.1.0}.
L'URL del repository è configurabile tramite \texttt{git.cloneUrl} in \texttt{config.yaml}.

\paragraph{Parametri di tuning configurabili (\texttt{MetricsConfig})}

Tre parametri che influenzano il tempo di esecuzione e il consumo di memoria della
Milestone 1 erano hardcoded come costanti statiche in \texttt{CodeSmellsMetric.java}.
Sono stati spostati nella nuova sezione \texttt{metrics} di \texttt{config.yaml}:

\begin{table}[h!]
\centering
\small
\begin{tabular}{@{}p{0.32\textwidth}p{0.3\textwidth}p{0.3\textwidth}@{}}
\toprule
\textbf{Parametro} & \textbf{Prima} & \textbf{Dopo} \\
\midrule
Frazione di snapshot processati &
argomento hardcoded passato a \texttt{computeAll()} &
\texttt{metrics.snapshotPercentage} in \texttt{config.yaml} \\
\addlinespace
File analizzati per invocazione PMD &
\texttt{static final int BATCH\_SIZE = 100} &
\texttt{metrics.pmd.batchSize} \\
\addlinespace
Frazione core assegnata a PMD &
\texttt{static final double CPU\_FRACTION = 0.5} &
\texttt{metrics.pmd.cpuFraction} \\
\bottomrule
\end{tabular}
\caption{Parametri di tuning spostati da costanti hardcoded a \texttt{config.yaml}.}
\end{table}

È stata aggiunta la classe \texttt{MetricsConfig} (con inner class \texttt{PmdConfig}),
mappata da SnakeYAML come tutte le altre sezioni di \texttt{AppConfig}. Il parametro
\texttt{snapshotPercentage} è particolarmente utile per i run di debug: impostarlo a
\texttt{0.1} processa solo il 10\% degli snapshot, riducendo il tempo di esecuzione da
ore a minuti.

\paragraph{Riorganizzazione delle cartelle di output}

Tutti gli output della Milestone 1 erano in cartelle con nomi generici
(\texttt{originalResult}, \texttt{enrichedResult}, ecc.) direttamente sotto
\texttt{output/}. Le cartelle sono state rinominate con nomi descrittivi e prefissi
numerici che riflettono l'ordine di produzione, e spostate sotto
\texttt{output/milestone1/}. Il CSV della Milestone 2 va sotto \texttt{output/milestone2/}.

\begin{table}[h!]
\centering
\small
\begin{tabular}{@{}ll@{}}
\toprule
\textbf{Vecchio percorso} & \textbf{Nuovo percorso} \\
\midrule
\texttt{output/originalResult/}   & \texttt{output/milestone1/1\_raw\_jira\_data/} \\
\texttt{output/enrichedResult/}   & \texttt{output/milestone1/2\_enriched\_jira\_data/} \\
\texttt{output/checkedResult/}    & \texttt{output/milestone1/3\_consistency\_checked/} \\
\texttt{output/proportionResult/} & \texttt{output/milestone1/4\_proportion\_applied/} \\
\texttt{output/snapshotResult/}   & \texttt{output/milestone1/5\_snapshots/} \\
\texttt{output/code/}             & \texttt{output/milestone1/6\_extracted\_source/} \\
\texttt{output/classifierResult/} & \texttt{output/milestone2/} \\
\bottomrule
\end{tabular}
\caption{Riorganizzazione delle cartelle di output. I prefissi numerici rendono leggibile
l'ordine di produzione dei dati lungo la pipeline.}
\end{table}

I path sono configurabili in \texttt{config.yaml} (campi \texttt{csv.outputDir},
\texttt{git.codeOutputDir}, \texttt{classifier.inputCsv}, \texttt{classifier.outputDir}).
Le directory vengono create automaticamente alla prima esecuzione se non esistono.
